\documentclass[11pt]{article}

\usepackage[a4paper,margin=2.5cm]{geometry}
\usepackage{amsmath,amssymb,mathtools}
\usepackage{microtype}
\usepackage[hidelinks]{hyperref}

\title{Relational Dynamic State Space\\
\large Algebraic Time, Hilbert--Moir\'e Geometry, Relational Attachment, and Page Attractors}
\author{Adrian Lipa}
\date{29 September 2026}

\newcommand{\AT}{\mathfrak{A}_{T}}
\newcommand{\HHM}{\mathcal{H}_{\mathrm{HM}}}
\newcommand{\Nrel}{\mathcal{N}_{\mathrm{rel}}}
\newcommand{\Apage}{\mathcal{A}_{\mathrm{page}}}

\begin{document}
\maketitle

\begin{abstract}
Relational Dynamic State Space (RDSS) is introduced as a framework in which
time and space are not postulated as coordinates of one primitive manifold.
Time is represented by an algebra of ordered transformations, while space is
represented by a Hilbert--Moir\'e state carrier equipped with an explicit
lattice/topology layer. Relational semantics produces a scalar attachment
field that deforms the metric of the carrier. Stable configurations of the
coupled state--metric dynamics are called page attractors. The construction
keeps arithmetic truth separate from geometric representation: an expression
such as \(2+2=4\) retains its algebraic value while its relational-semantic
embedding may induce a non-trivial scalar deformation. No identification
with physical mass or gravity is assumed at this foundational stage.
\end{abstract}

\section{Foundational separation}

RDSS starts from a strict separation:
\[
\boxed{\text{time}=\text{algebraic dynamics}},
\qquad
\boxed{\text{space}=\text{geometric state carrier}}.
\]
Let \(\AT\) denote the algebra of admissible ordered transformations and let
\(\HHM\) denote the Hilbert--Moir\'e carrier. Their coupling is expressed by
a representation
\[
\rho:\AT\longrightarrow \operatorname{End}(\HHM).
\]
A discrete orbit therefore has the form
\[
\Psi_{n+1}=\rho(a_n)\Psi_n,
\qquad a_n\in\AT,
\]
without requiring a primitive external real-valued time coordinate.

\section{Hilbert--Moir\'e space}

The spatial layer is written schematically as
\[
\HHM=\mathcal{H}(\Lambda_{\mathrm{HO}},\mathcal F),
\]
where \(\Lambda_{\mathrm{HO}}\) is a hyperoctagonal-lattice structure and
\(\mathcal F\) is a local state fibre. At the current stage, RDSS treats the
precise hyperoctagonal realization as a replaceable topology contract rather
than silently canonizing a single combinatorial model.

A useful factorized representation is
\[
\HHM \simeq \ell^2(\Lambda_{\mathrm{HO}})
\otimes \mathcal H_{\mathrm{local}},
\]
augmented by a Moir\'e/refinement structure.

\section{Relational attachment}

Let \(\mathcal R\) denote relational structure, \(C\) context, and \(\Psi\)
the current state. The relational attachment is the map
\[
\boxed{
\Nrel:(\Psi,\mathcal R,C)\longmapsto \chi
}
\]
with scalar output \(\chi\in\mathbb R\).

The reference implementation uses an explicit feature vector \(f_k\) and a
replaceable linear functional
\[
\chi=b+\sum_k w_k f_k.
\]
This is an implementation baseline, not a claim that the physical or
canonical semantic functional must be linear.

\section{Semantic scalar deformation}

Given a base metric \(g^{(0)}\), RDSS v0.1 uses the conformal deformation
\[
\boxed{
g[\chi]=e^{2\chi}g^{(0)}.
}
\]
This choice is explicit, differentiable, positive for finite real \(\chi\),
and easy to falsify or replace.

The critical invariant is that arithmetic evaluation and geometric
representation remain separate channels. For
\[
E=(2+2)
\]
the value channel gives
\[
\operatorname{value}(E)=4,
\]
while the relational-semantic channel may give
\[
\Nrel(E,C)=\chi_E
\]
and consequently
\[
g^{(0)}\longmapsto e^{2\chi_E}g^{(0)}.
\]
Thus RDSS does not alter arithmetic truth. It allows the representation of a
relation to modify the geometry of the state carrier.

\section{Coupled dynamic state}

A coupled state is represented by
\[
X_n=(\Psi_n,g_n).
\]
One RDSS step is
\[
\Psi_{n+1}=\rho(a_n)\Psi_n,
\]
followed by
\[
\chi_{n+1}=\Nrel(\Psi_{n+1},\mathcal R_{n+1},C_{n+1}),
\]
and
\[
g_{n+1}=\mathcal D_{\chi_{n+1}}(g^{(0)}).
\]
In the initial reference kernel
\[
\mathcal D_{\chi}(g^{(0)})=e^{2\chi}g^{(0)}.
\]

This yields the structural loop
\[
\boxed{
\text{algebraic action}
\rightarrow
\text{state}
\rightarrow
\text{relational semantics}
\rightarrow
\chi
\rightarrow
\text{metric deformation}.
}
\]

Future versions may allow \(g_{n+1}\) itself to modify the representation
\(\rho\), producing a fully bidirectional feedback loop.

\section{Page attractor}

A page attractor is a stable configuration of the coupled update.
For an update map
\[
\mathcal F:(\Psi_n,g_n)\mapsto(\Psi_{n+1},g_{n+1}),
\]
a fixed-point page attractor satisfies
\[
\boxed{
\mathcal F(\Psi_\star,g_\star)=(\Psi_\star,g_\star).
}
\]
The v0.1 computational definition declares convergence when both
\[
\|\Psi_{n+1}-\Psi_n\|_2\le\varepsilon
\]
and
\[
\|g_{n+1}-g_n\|_F\le\varepsilon.
\]
Periodic, quasiperiodic, and chaotic attractors are deliberately deferred.


\section{Hilbert primacy and derived representations}

RDSS has one canonical carrier:
\[
\boxed{\HHM.}
\]
No Bloch sphere, K\"ahler chart, Minkowski chart, Euclidean chart or other
geometry is introduced as a second foundational space. Such geometries are
derived representations of states already living in \(\HHM\).

For each admissible representation label \(\alpha\), let
\[
\boxed{
\Pi_\alpha:\HHM\longrightarrow \mathcal X_\alpha
}
\]
be a projection, chart, quotient, embedding coordinate map, or other explicit
representation morphism. The symbol \(\mathcal X_\alpha\) denotes the image
or coordinate realization selected by \(\Pi_\alpha\); it carries no
independent ontological status in RDSS.

A metric or quadratic form belongs to a representation only after such a
choice. If \(G_\alpha\) is the corresponding form, then for two Hilbert--Moir\'e
states \(\Psi,\Phi\) one may define
\[
Q_\alpha(\Psi,\Phi)
=
\left\langle
\Pi_\alpha(\Psi-\Phi),
G_\alpha\Pi_\alpha(\Psi-\Phi)
\right\rangle.
\]
Thus a Lorentzian/Minkowski form, a Bloch representation, or a K\"ahler/
Fubini--Study description is a choice of \((\Pi_\alpha,G_\alpha)\) inside the
Hilbert framework, not a replacement for it.

The informational--holonomic coordinate is likewise internal:
\[
\boxed{
\zeta(\Psi)=I(\Psi)+i\Theta(\Psi).
}
\]
It is a complex degree of freedom or functional of the Hilbert state. It is
not postulated as an external spacetime axis.

\section{Holonomy inside Hilbert--Moir\'e space}

Let \(\Gamma\) be a trajectory in the Hilbert--Moir\'e state space (or its
projective realization when global phase is quotiented). Relational transport
is described by a connection \(\mathcal A_{\rm rel}\) and holonomy
\[
\boxed{
U_\Gamma
=
\mathcal P\exp\!\left(
i\int_\Gamma \mathcal A_{\rm rel}
\right).
}
\]
A representation may close even when the complete Hilbert state does not.
For some representation \(\Pi_\alpha\), it is possible that
\[
\Pi_\alpha(\Psi_{\rm final})
=
\Pi_\alpha(\Psi_{\rm initial}),
\]
while simultaneously
\[
\boxed{
\Psi_{\rm final}
=
U_\Gamma\Psi_{\rm initial}
\neq
\Psi_{\rm initial}.
}
\]
This is the precise RDSS version of returning to the same represented point
while retaining path-dependent information. Closure therefore has to be
qualified by representation:
\[
\boxed{
\text{representation closure}
\not\Rightarrow
\text{Hilbert-state closure}.
}
\]

For a scalar \(U(1)\) reference transport,
\[
U_\Gamma
=
\exp\!\left(
i\sum_{e\in\Gamma}\phi_e
\right),
\]
with winding and residual
\[
W_\Gamma
=
\operatorname{round}\!\left(
\frac{1}{2\pi}\sum_{e\in\Gamma}\phi_e
\right),
\qquad
\varepsilon_\Gamma
=
\operatorname{wrap}_{2\pi}\!\left(
\sum_{e\in\Gamma}\phi_e
\right).
\]
Higher-dimensional and non-Abelian transports remain admissible extensions.

\section{Representation-specific geometries are derived, not fundamental}

The time algebra continues to act on the canonical carrier,
\[
\rho:\AT\longrightarrow\operatorname{End}(\HHM).
\]
Only after a representation \(\Pi_\alpha\) is selected does one obtain a
representation-specific geometric description. Schematically,
\[
\boxed{
\HHM
\xrightarrow{\ \Pi_\alpha\ }
\mathcal X_\alpha,
\qquad
\Psi\mapsto x_\alpha.
}
\]
A Lorentzian representation may therefore use a Minkowski form, while a
qubit sector may be represented projectively on a Bloch sphere with its
K\"ahler/Fubini--Study geometry. These are examples of derived descriptions
of Hilbert states. RDSS does not promote any one of them to the universal
carrier.

Semantic dependence may act directly on Hilbert operators, representation
maps, representation metrics, or the relational connection. A generic form
is
\[
\chi=\Nrel(\Psi,\mathcal R,C),
\]
followed by one or more deformations such as
\[
G_\alpha\longmapsto G_\alpha[\chi],
\qquad
\mathcal A_{\rm rel}\longmapsto\mathcal A_{\rm rel}[\chi].
\]
The representation can therefore retain the same symbolic form while the
relational/holonomic content of the transformation changes.

\section{Core object}

The minimal RDSS object is therefore
\[
\boxed{
\left(
\HHM,\,
\AT,\,
\rho,\,
\Nrel,\,
\mathcal D,\,
\Apage
\right).
}
Equivalently, the evolving geometry can be displayed as
\[
\boxed{
\AT
\curvearrowright
\left[
\HHM,
g\!\left(
\Nrel[\operatorname{Sem}(\mathcal R,\Psi)]
\right)
\right].
}
\]

\section{Epistemic boundary}

The following statements belong to the mathematical framework at this stage:
\begin{itemize}
  \item time dynamics and spatial carrier are represented separately;
  \item an explicit relational attachment can produce a scalar \(\chi\);
  \item \(\chi\) can explicitly deform a metric;
  \item coupled state--metric updates can possess attractors.
\end{itemize}

The following identifications are \emph{not} assumed:
\begin{itemize}
  \item page attractor \(=\) physical mass;
  \item relational attachment/deformation \(=\) physical gravity;
  \item the current linear semantic functional is canonical;
  \item the current conformal deformation law is uniquely forced.
\end{itemize}

Those interpretations require separate derivation and validation gates.

\section{Immediate research programme}

The next mathematical gates are:
\begin{enumerate}
  \item define and test candidate hyperoctagonal adjacency/topology models;
  \item formalize the generators and composition law of \(\AT\);
  \item derive semantic features from relational structure rather than insert
        them by hand;
  \item classify fixed, periodic, and higher-order page attractors;
  \item determine invariants under semantic deformation;
  \item test whether geometry can feed back into \(\rho\) without
        inconsistency;
  \item only then test candidate physical interpretations.
\end{enumerate}

\end{document}
